 % \documentclass[10pt,respuestas,a4]{examenEPN}
\documentclass[10pt,a4]{examenEPN0}
\usepackage[T1]{fontenc}
\examen{Prueba 01}
\usepackage{comandosEPN}

\usepackage{empheq}

\newcommand{\itemcp}{\stepcounter{enumi}\item[ \textcolor{red}{\theenumi.$^*$}] }

\newcommand{\iitemcp}{\stepcounter{enumii}\item[ \textcolor{red}{\theenumii)$^*$}] }


\materia{Álgebra Lineal}
%\autor{Departamento de Formación Básica}
\autor{GR \hspace{1cm}.}
%\fecha{Semestre 2026-A}
\fecha{Nombre: \hspace{9cm} Cédula:}
\logouno[1.6cm]{Logo_EPN.eps}
\logodos[2.6cm]{Logo_DFB.eps}

\newcommand{\la}{\lambda}

%Datos
\newcommand{\datosAlgebra}[2]{
\EnExamen
{
\vspace{5mm}
    \noindent
    Nombres: \hrulefill\
    Carrera: \rule{4.5cm}{0.4pt}\
    \ifthenelse{\equal{#1}{}}%
    {Paralelo: \underline{\parbox{1.8cm}{\parbox{1cm}{}}}\\[3mm]}
    {Paralelo: \underline{\parbox{1.8cm}{\centering #1}}\\[3mm]}
    C\'edula: \rule{5cm}{0.4pt}\
    Firma: \hrulefill\
    Fecha: \underline{\parbox{2.1cm}{\centering #2}}
}}


% \newcommand{\indicacionesPrueba}{%
% \EnExamen{
% \begin{indicaciones}
%     \item La prueba tiene una duraci\'on  de \textbf{120 minutos.}
%     \item No se permite la utilizaci\'on de dispositivos electr\'onicos (celular, tablets, calculadora, etc.)
%     \item Leer detenidamente los enunciados, \textbf{justificar la respuesta para acceder al puntaje total}.
%     \item La prueba es completamente individual. Los estudiantes deberán dejar sus teléfonos celulares guardados y no podrán tenerlos consigo (por ejemplo, en los bolsillos de su ropa) durante toda la evaluación. En caso de detectarse algún acto de deshonestidad académica o el incumplimiento de esta norma, no podrán seguir rindiendo la prueba y su calificación será \textbf{cero}.
% \end{indicaciones}
% \pagestyle{empty}
% }}

\newcommand{\indicacionesPrueba}{%
\EnExamen{
\begin{indicaciones}
    \item La prueba tiene una duraci\'on  de \textbf{60 minutos.}
    \item Leer detenidamente los enunciados, \textbf{justificar la respuesta} para acceder al puntaje total.
    \item No se permite la utilizaci\'on de dispositivos electr\'onicos (celular, tablets, calculadora, etc.)
    \item La prueba es completamente individual
    \item No se permite portar teléfonos celulares, relojes inteligentes ni otros dispositivos electrónicos durante el examen, aun cuando permanezcan en los bolsillos. Estos dispositivos deberán guardarse en la mochila antes del inicio de la evaluación y esta deberá mantenerse cerrada durante todo el examen.
    \item He leído las indicaciones y acepto que el incumplimiento de estas se penaliza con \textbf{cero}. Firma:...................................

    
\end{indicaciones}
\pagestyle{empty}
}}

\newcommand{\indicacionesExamen}{
\EnExamen{\begin{indicaciones}
    \item La evaluación tiene una duraci\'on  de \textbf{120 minutos.}
    \item No se permite la utilizaci\'on de dispositivos electr\'onicos (celular, tablets, calculadora, etc.)
    %\item Leer detenidamente los enunciados, justificar la respuesta para acceder al puntaje total.
    \item El examen es completamente individual. En caso de detectarse alg\'un acto de deshonestidad acad\'emica no podr\'a seguir rindiendo la prueba y su calificaci\'on ser\'a \textbf{cero}.
\end{indicaciones}
}\pagestyle{empty}
}

%% Imprimir las una hoja A4, dos veces en A5
\newcommand{\imprimirDoble}
{%
\usepackage{pgfmorepages}
\pgfmorepagesloadextralayouts
\pgfpagesuselayout{repeated 2-up}[a4paper,landscape]
\thispagestyle{empty}
}


%Este sí se puede comentar
\newenvironment{solucion}{\color{blue}\proof[Solución]}{\endproof\color{black}}




% \renewcommand{\@autor}{
% ---\\
% \phantom{.} 
% ...}
% \makeatother

%\imprimirDoble
\begin{document}
\encabezado
\datosAlgebra{GR5E}{22/10/2026}
%Poner el grupo y la fecha, o dejar las llaves vacías.
\indicacionesPrueba
\begin{preguntas}

\item Sea $\alpha\in\R$. Halle el rango de la matriz en función de $\alpha$:
\[
A = \begin{pmatrix} 
1 & -1 & 2 \\ 
2 & -2 & 5 \\ 
-1 & \alpha & -1 
\end{pmatrix} 
\]
\puntaje{5}
% Espacio de respuesta hasta el final de la página.
\par\vfill\newpage


%%%%%%%%%%%%%%%%%%%%%%%%%%%%%%%%%%


\item Demuestre o refute:

    %\item Sean que $A, B \in \mathbb{R}^{n \times n}$ matrices simétricas entonces $AB - BA$ es antisimétrica
  Sea $A \in \mathbb{R}^{n \times n}$  una matriz idempotente entonces la matriz $B = 2A - I_n$ es involutiva (es decir, $B^2 = I_n$).
\puntaje{5}
% Espacio de respuesta hasta el final de la página.
\par\vfill\newpage






%%%%%%%%%%%%%%%%%%%%%%%%%%%%%%%%%

\item Considere el siguiente sistema de ecuaciones lineales en las variables $x, y, z$: 
\[
\begin{cases} 
x - y + z = 0 \\ 
kx + 2y  -2z = -k \\ 
 y + kz = 1 
\end{cases}
\]

\begin{enumerate}
 \item  Escriba la matriz ampliada $(A \mid b)$ asociada al sistema y aplique, consecutivamente, las siguientes operaciones elementales por fila  \puntaje{1} \[-kF_1 + F_2\rightarrow F_2;\quad F_2\leftrightarrow F_3 \quad  \text{y} \quad  -(2+k)F_2 + F_3 \rightarrow F_3\]


\item  Determine, utilizando el Teorema de Rouché-Frobenius, los valores de $k\in \mathbb{R}$ para los cuales el sistema tiene solución única, infinitas soluciones o no tiene solución. Además, si $k=-1$, obtenga el conjunto solución.  \puntaje{4}
\end{enumerate}
% Espacio de respuesta hasta el final de la página.
\par\vfill\newpage


%%%%%%%%%%%%%%%%%%%%%%%%%%%%%%%%%%%%%%%

\item Una planta de manufactura automatizada utiliza tres máquinas de control numérico computarizado (Maquinaria $X$, $Y$, $Z$) para fabricar tres piezas críticas de un motor aeronáutico. El tiempo de operación (en minutos) que requiere cada pieza en cada máquina es el siguiente:

\begin{itemize}
    \item La \textbf{Pieza 1} requiere 1 minuto en la máquina $X$, 2 minutos en la $Y$, y 3 minutos en la $Z$.
    \item La \textbf{Pieza 2} requiere 3 minutos en la máquina $X$, 1 minuto en la $Y$, y 2 minutos en la $Z$.
    \item La \textbf{Pieza 3} requiere 2 minutos en la máquina $X$, 3 minutos en la $Y$, y 1 minuto en la $Z$.
\end{itemize}

Si al final de la jornada las máquinas registraron un tiempo de operación total de \textbf{46 minutos la Máquina X}, \textbf{36 minutos la Máquina Y} y \textbf{38 minutos la Máquina Z}. Plantee el sistema de ecuaciones que determine cuántas unidades de cada pieza se fabricaron durante el turno. \puntaje{5}


%%%%%%%%%%%%%%%%%%%%%%%%%%%%%%%%%%%%%%%%%

    
% El espacio restante de la última página queda disponible para responder.
\par\vfill
\end{preguntas}
\final
\end{document}
